\chapter{Введение}
\label{ch:intro}

Жители городов, находящихся рядом с заводами, на которых используются опасные для жизни человека газы, могут попасть в зону поражения, если на предприятии произойдёт авария. В такой ситуации людям нужно как можно скорее получить информацию о происшествии и определить план действий для избежания угрозы. 

Есть несколько способов получения информации. Первый подразумевает использование карманного газоанализатора, для контроля концентрации опасного газа в воздухе. Однако, когда уровень концентрации превысит норму, эвакуироваться может быть поздно. Также, носить газоанализатор с собой и контроллировать концентрацию вручную - неудобно.

Есть другой вариант. Человек может узнать о происшествии из СМИ и рассчитать скорость распространения газа в воздухе, зная коэффициент диффузии. Этот способ позволяет быстро оценить обстановку и составить план эвакуации.

Целью работы являлась проверка метода определения коэффициента диффузии на примере гелия.
